% ======================================================================
%  40_layer3_terms.tex — 细度层：术语、符号与常见误解
%  依赖 preamble.tex；独立自检目录 build/check_terms/
% ======================================================================

\section{术语、符号与常见误解}

读原文时真正卡住人的，往往不是某个公式，而是两件更基本的事：一个词在作者笔下到底指什么，一个记号第一次在哪儿开始承担技术含义。这一章把这两件事先钉死，再逐条拆掉最容易走偏的直觉。

三张表对应三个问题：\emph{这个词在原文里指什么}（\S\ref{sec:terms-glossary}）、\emph{这个记号在哪里第一次真正被使用}（\S\ref{sec:terms-symbols}）、\emph{我凭常识以为的意思为什么不对}（\S\ref{sec:terms-misconceptions}）。三张表都不是百科条目，而是为读这篇论文服务的工具：每一条都尽量挂回原文的某一页，方便你翻回去核对。

\begin{kcard}{本层的三条约定}
\begin{enumerate}[label=\arabic*.,leftmargin=1.7em,itemsep=2pt]
  \item \textbf{首现只记技术首现。} \S1 是路线图，很多术语在那里只是被点名（例如 A-模型、几何 Langlands、Fourier--Mukai 变换），本表记的是它在正文中第一次被定义或被实质使用的位置。两处都重要时写成「\S1 p.2 / \S5.1 p.65」。
  \item \textbf{英文列保留原文拼写。} 连字符、撇号一律照抄：'t Hooft 不写成 t Hooft；S-对偶、T-对偶带连字符；原文写 $N=4$ 而不写 $\mathcal{N}=4$；\emph{Hecke eigensheaf} 的「eigen」是「本征」不是「特征值」。
  \item \textbf{「易混」列只写致命的混淆。} 同义词辨析、译名之争不进这一列。进这一列的都是「按字面理解就会推出错误结论」的那种，例如把扭化当成简化理论、把 eigenbrane 当成 Hilbert 空间里的本征态。
\end{enumerate}
\end{kcard}

\subsection{术语表}\label{sec:terms-glossary}

全表按主题分八组：A 对偶与群论、B 超对称与扭化、C Hitchin 模空间与超凯勒几何、D 广义复几何与膜、E 线算子与特征膜、F 单极子与 Hecke 几何、G Langlands 与 D-模、H 原文只引用而未真正使用的背景术语（标 $\dagger$ 者全文一次也没有以该词形出现，表中给出原文的实际说法）。\anchor{src p.2}

\begingroup
\footnotesize
\setlength{\tabcolsep}{5pt}
\renewcommand{\arraystretch}{1.12}
\begin{longtable}{@{}>{\raggedright\arraybackslash}p{2.45cm}>{\raggedright\arraybackslash}p{1.78cm}>{\raggedright\arraybackslash}p{5.35cm}>{\raggedright\arraybackslash}p{1.78cm}>{\raggedright\arraybackslash}p{2.95cm}@{}}
\toprule
\tabhead{英文} & \tabhead{中文} & \tabhead{一句话定义} & \tabhead{首现} & \tabhead{易混淆点} \\
\midrule
\endfirsthead
\toprule
\tabhead{英文} & \tabhead{中文} & \tabhead{一句话定义} & \tabhead{首现} & \tabhead{易混淆点} \\
\midrule
\endhead
\bottomrule
\endlastfoot

\multicolumn{5}{@{}l}{\tabhead{A\quad 对偶、群论与耦合常数（\S1--\S2）}}\\
\addlinespace[1pt]

electric--magnetic duality & 电-磁对偶 & 交换电、磁荷并反演耦合的理论对称 & \S1 p.2 & 不只是「把 $\bm E$ 换成 $\bm B$」，必须同时换群 \\
\addlinespace[1pt]
Montonen--Olive duality & Montonen--Olive 对偶 & 同一超对称规范理论在 $g\leftrightarrow 1/g$ 下等价的猜想 & \S1 p.2 & 原猜想只含 $\ZZ_2$；加入 $\theta$ 后推广为 $SL(2,\ZZ)$ \\
\addlinespace[1pt]
S-duality & S-对偶 & $\tau\mapsto-1/n_g\tau$、$G\leftrightarrow\Ldual$ 的量子对称 & \S1 p.4 / \S2.2 p.11 & 只在具有足够超对称的理论里成立 \\
\addlinespace[1pt]
T-duality & T-对偶 & 2d sigma 模型沿环面反演半径的对偶 & \S1 p.3 & 四维没有 T-对偶；它是紧化后的二维现象 \\
\addlinespace[1pt]
GNO dual group & GNO 双群 & 由电、磁荷取值格的对偶确定的群 & \S1 p.3 & 与 Langlands 对偶群是同一个群，只是出发点在格 \\
\addlinespace[1pt]
Langlands dual group & Langlands 对偶群 & 根与余根数据互换所得的群 $\Ldual$ & \S1 p.3 & $SU(N)$ 的对偶是 $PSU(N)$，不是 $SU(N)$ 自己 \\
\addlinespace[1pt]
simply-laced & 单连通型 & 所有根等长（$A,D,E$ 型） & \S2.2 p.11 & 此时 $S:\tau\mapsto-1/\tau$；否则 $n_g=2,3$ \\
\addlinespace[1pt]
duality group & 对偶群 $\Gamma$ & $T$ 与 $S$ 生成的 $\Gamma\subset SL(2,\RR)$ & \S2.2 p.11 & 不是规范群 $G$，是作用在 $\tau$ 上的群 \\
\addlinespace[1pt]
theta angle & $\theta$ 角 & $F\wedge F$ 拓扑项的系数 & \S1 p.3 / \S2.2 p.11 & 只是 $\tau$ 的实部；不是独立的耦合 \\
\addlinespace[1pt]
complexified coupling & 复耦合 $\tau$ & $\tau=\theta/2\pi+4\pi\ii/e^2$ & \S2.2 p.11 & 与 canonical parameter $\Psi$ 不是一个东西 \\
\addlinespace[1pt]
weight lattice & 权格 & 电荷取值的格 & \S2.2 p.13 & 磁荷取的是对偶群的权格，别写反 \\
\addlinespace[1pt]
coweight lattice & 余权格 & 磁荷取值的格 & \S2.2 p.13 & 模掉 Weyl 群后才唯一 \\
\addlinespace[1pt]
GNO line & GNO 线 & 按 $\Ldual$ 的表示给 't Hooft 算子贴的标签 & \S6.2 p.84 & 不是几何意义上的线，是表示的标签 \\
\addlinespace[1pt]
Hecke group & Hecke 群 & $G_2,F_4$ 时完整对偶群所在的 $SL(2,\RR)$ 子群 & \S2.2 p.11 & 不要与 $\ZZ_2$ 型的经典对偶混为一谈 \\

\addlinespace[4pt]
\multicolumn{5}{@{}l}{\tabhead{B\quad 超对称、扭化与拓扑场论（\S2--\S3、\S6.4）}}\\
\addlinespace[1pt]

$N=4$ supersymmetric Yang--Mills & $N=4$ 超对称杨--米尔斯 & 四维最大超对称的规范理论 & \S1 p.3 & 不是 $N=4$ 超引力，这里没有引力 \\
\addlinespace[1pt]
supercharge & 超荷 & 超对称代数中把玻色子变为费米子的生成元 & \S2.1 p.9 & 与共形超荷不是一回事 \\
\addlinespace[1pt]
scalar supercharge & 标量超荷 & 扭化后在转动下为标量且 $Q^2=0$ 的超荷 & \S3.1 p.18 & 原文多称 topological supercharge；此原串未出现 \\
\addlinespace[1pt]
BRST operator & BRST 算子 & 扭化理论中 $Q$ 的别名 & \S6.4 p.97 & 不是规范固定意义下的 BRST \\
\addlinespace[1pt]
cohomology of $Q$ & $Q$-上同调 & 只留 $Q$-闭、模去 $Q$-恰当的态与算子 & \S3.1 p.18 & 漏掉「模去恰当」就只剩不变子空间 \\
\addlinespace[1pt]
R-symmetry & R-对称 & $N=4$ 的内禀 $SU(4)_R$ 对称 & \S2.1 p.10 & 扭化正是把它的子代数混入 Lorentz 群 \\
\addlinespace[1pt]
twisting & 扭化 & 经同态 $\kappa$ 把自旋群换成对角子群、重定义自旋 & \S3.1 p.17 & 不是「简化理论」；作用变了，理论还是同一个 \\
\addlinespace[1pt]
GL twist & GL 扭化 & 使 $S^+$ 含不变向量、从而得到标量 $Q$ 的扭化 & \S3.1 p.17 & 与 Donaldson--Witten 扭化不同 \\
\addlinespace[1pt]
topological field theory & 拓扑场论 & 能量--动量张量 $Q$-恰当、只依赖拓扑的理论 & \S1 p.2 & 原文只说「松散地说是 TQFT」 \\
\addlinespace[1pt]
instanton number & 瞬子数 & $(1/8\pi^2)\int\Tr F\wedge F$，即 $c_2$ & \S1 p.3 / \S2.2 p.12 & 方向反转时变号（故 $T:\Psi\mapsto-\Psi$） \\
\addlinespace[1pt]
BPS condition & BPS 条件 & 质量等于中心荷、保住部分超对称的条件 & \S5.4 p.73 & 不是「质量为零」；半 BPS 只保一半 \\
\addlinespace[1pt]
central charge & 中心荷 & 超对称代数的中心元，给出 BPS 质量 & \S2.2 p.14 & 与二维 CFT 的中心荷无关 \\
\addlinespace[1pt]
topological supercharge & 拓扑超荷 & 标量超荷 $Q$ 的另一种叫法 & \S6.4 p.102 & 同一对象有三个名字：标量/拓扑/BRST \\
\addlinespace[1pt]
Coulomb branch & Coulomb 分支 & 标量期望值非零、规范群被 Higgs 化的分支 & \S2.2 p.14 & 本文的 $\MH$ 不是 Coulomb 分支 \\

\addlinespace[4pt]
\multicolumn{5}{@{}l}{\tabhead{C\quad Hitchin 模空间与超凯勒几何（\S3--\S4）}}\\
\addlinespace[1pt]

Hitchin moduli space & Hitchin 模空间 & 稳定 Higgs 丛的模空间 $\MH$ & \S4.1 p.45 & 不是 $G$-丛模空间 $\mathcal{M}$ \\
\addlinespace[1pt]
Higgs bundle & Higgs 丛 & 带全纯 $1$-形式 $\varphi$ 的丛 $(E,\varphi)$ & \S4.2 p.49 & 不是 Higgs 场的真空模空间（那是 $\MH$） \\
\addlinespace[1pt]
Higgs field & Higgs 场 & $\varphi\in H^0(C,K_C\otimes\mathrm{ad}E)$ & \S4.2 p.49 & 与四维理论里的标量不是同一个对象 \\
\addlinespace[1pt]
Hitchin pair & Hitchin 对 & 满足 $\bar D\varphi=0$ 且稳定的 $(E,\varphi)$ & \S4.2 p.49 & 稳定性条件不可省 \\
\addlinespace[1pt]
stability & 稳定性 & 对 $\varphi$-不变子丛的斜率条件 & \S4.2 p.50 & 忘掉它只得到拟模空间 \\
\addlinespace[1pt]
hyper-Kähler manifold & 超凯勒流形 & 对三个复结构都是 Kähler 的流形 & \S4.1 p.42 & 复维必须为偶数，且 $I,J,K$ 满足四元数关系 \\
\addlinespace[1pt]
hyper-Kähler quotient & 超凯勒商 & 三个 moment map 同时取零再取商 & \S4.1 p.40 & 不是普通 Kähler 商 \\
\addlinespace[1pt]
complex structures $I,J,K$ & 复结构三元组 & $I^2=J^2=K^2=IJK=-1$ 的三元组 & \S4.1 p.42 & $J$ 与 $-J$ 不依赖 $C$ 的复结构 \\
\addlinespace[1pt]
twistor $\PP^1$ & 复结构 $\PP^1_g$ & 复结构族 $I_w$ 的参数空间，$w\in\PP^1$ & \S4.1 p.43 & 不是靶空间的一族度量 \\
\addlinespace[1pt]
symplectic form & 辛形式 & 闭非退化 $2$-形式 $\omega_I,\omega_J,\omega_K$ & \S4.1 p.44 & $\omega_J,\omega_K$ 恰当，$\omega_I$ 拓扑非平凡 \\
\addlinespace[1pt]
moment map & moment map $\mu$ & 由 $\iota_{V(\epsilon)}\omega=\delta\mu$ 定义的映射 & \S4.1 p.45 & 三个 moment map 同为零才是 Hitchin 方程 \\
\addlinespace[1pt]
Hitchin fibration & Hitchin 纤维化 & $\pi:\MH\to B$，只保留 Casimir 多项式 & \S4.3 p.53 & 与「忘掉 $\varphi$ 的第二纤维化」不是同一个映射 \\
\addlinespace[1pt]
integrable system & 可积系统 & $B$ 上线性函数在 Poisson 括号下交换 & \S4.3 p.54 & 可积不等于能显式写出解 \\
\addlinespace[1pt]
Casimir basis & Hitchin 基 $B$ & $\bigoplus_i H^0(C,K_C^{d_i})$ & \S4.3 p.53 & 维数为 $\frac12\dim\MH$ \\
\addlinespace[1pt]
spectral cover & 谱覆盖 & 由 $\mathrm{det}(\varphi-\lambda)$ 定义的 $B$ 上的覆盖 & \S4.3 p.54（原文声明不用） & 本文完全没有用这个技术 \\

\addlinespace[4pt]
\multicolumn{5}{@{}l}{\tabhead{D\quad 广义复几何、sigma 模型与膜（\S5）}}\\
\addlinespace[1pt]

sigma model & sigma 模型 & 以映射 $\Phi:\Sigma\to X$ 为基本场的场论 & \S1 p.2 & 此处指二维超对称拓扑版本 \\
\addlinespace[1pt]
A-model & A-模型 & 只依赖复化 Kähler 类 $[B+\ii\omega]$ 的拓扑 sigma 模型 & \S1 p.2 / \S5.1 p.65 & 与 A-膜是两个层次的概念 \\
\addlinespace[1pt]
B-model & B-模型 & 只依赖靶空间复结构的拓扑 sigma 模型 & \S1 p.4 / \S5.1 p.65 & 与 B-膜不是同一件事 \\
\addlinespace[1pt]
mirror symmetry & 镜像对称 & $X$ 与 $\tilde X$ 的 A/B 模型互换而等价的对称 & \S1 p.2 / \S5.5 p.74 & 不是两个 Calabi--Yau 流形的几何等同 \\
\addlinespace[1pt]
generalized complex geometry & 广义复几何 & $T\oplus T^*$ 上带自然度量的复结构理论 & \S5.2 p.65 & 不是复几何与辛几何的并集 \\
\addlinespace[1pt]
generalized complex structure & 广义复结构 & $I^2=-1$、保度量、可积的 $T\oplus T^*$ 结构 & \S5.2 p.66 & $I_{\hat J}$ 与 $I_\omega$ 只是两个极端例子 \\
\addlinespace[1pt]
B-field & B-场 & 与 $\omega$ 并列的 $2$-形式，作用量含 $\int\Phi^*(\omega-\ii B)$ & \S4.1 p.46 & 不是规范场；整体上非平凡 \\
\addlinespace[1pt]
coisotropic & 余各向同性 & Poisson 交换函数零点定义的子流形 & \S11.1 p.185 & 不是各向同性；$\dim Y\ge\frac12\dim X$ \\
\addlinespace[1pt]
A-brane & A-膜 & 边界落在 Lagrange 或余各向同性子流形上的膜 & \S5.2 p.69 & 不是 A-模型 \\
\addlinespace[1pt]
B-brane & B-膜 & 支撑在复子流形上、带全纯 Chan-Paton 丛的膜 & \S5.3 p.70 & 不是 B-模型 \\
\addlinespace[1pt]
brane & 膜 & 带边理论的边界条件，数据为支撑与 Chan-Paton & \S1 p.2 & 原文不特指 D-膜；除特例外不要加 D \\
\addlinespace[1pt]
boundary condition & 边界条件 & $\Sigma$ 边界上场的局部条件 & \S5.6 p.77 & 拓扑不变量要看它是否 $Q$-不变 \\
\addlinespace[1pt]
Chan-Paton bundle & Chan-Paton 丛 & 膜上的向量丛或层，承载规范量子数 & \S6.4 p.95 & 特征值就住在这个空间里 \\
\addlinespace[1pt]
zerobrane & 零膜 & 支撑在一点、Chan-Paton 为该点天穹层的膜 & \S5.5 p.74 & 并不平凡：它是电特征膜的原型 \\
\addlinespace[1pt]
type F brane & type F 膜 & 支撑在 Hitchin 纤维上、带平坦 $U(1)$ 丛的膜 & \S9 p.125 & 属 $(B,A,A)$ 型 \\
\addlinespace[1pt]
canonical coisotropic A-brane & 典范余各向同性 A-膜 & 秩 $1$、$(\omega^{-1}F)^2=-1$ 的余各向同性 A-膜 & \S1 p.5 / \S11.1 p.185 & 不是某个具体的 Lagrange 膜 \\
\addlinespace[1pt]
brane type & 膜的类型 & 对 $I,J,K$ 三个复结构分别标 A 或 B & \S5.6 p.77 & 三个字母对应三个复结构，不是三种膜 \\
\addlinespace[1pt]
SYZ picture & SYZ 图景 & 镜像对称由对偶环面纤维化实现 & \S5.5 p.77 & 原文只用它作启发，不当作证明 \\

\addlinespace[4pt]
\multicolumn{5}{@{}l}{\tabhead{E\quad 线算子、特征膜与 Hecke 算子（\S6、\S9）}}\\
\addlinespace[1pt]

Wilson operator & Wilson 算子 & 沿曲线取 holonomy 的迹，属 order operator & \S6.1 p.79 & 其超对称版本只在 $\Psi=\infty$ 存在 \\
\addlinespace[1pt]
't Hooft operator & 't Hooft 算子 & 规定经典场的奇异性再作路径积分，属 disorder operator & \S6.2 p.81 & 不是同一理论里 Wilson 的「对偶算符」 \\
\addlinespace[1pt]
line operator & 线算子 & 支撑在一维子流形上的算子 & \S6 p.79 & 开线与闭圈的算子代数不同 \\
\addlinespace[1pt]
loop operator & 圈算子 & 支撑在闭曲线上的线算子 & \S6.1 p.79 & 不要与局部算子混为一谈 \\
\addlinespace[1pt]
order operator & 有序算子 & 把经典量再量子化得到的算子 & \S6.2 p.81 & 与 disorder operator 相对 \\
\addlinespace[1pt]
disorder operator & 无序算子 & 规定奇异性后作路径积分得到的算子 & \S6.2 p.81 & 需要法丛定向 \\
\addlinespace[1pt]
eigenbrane & 特征膜 & 被线算子映到自身张量一个向量空间的膜 & \S6.4 p.95 & 本征值是向量空间或层，不是数 \\
\addlinespace[1pt]
electric eigenbrane & 电特征膜 & Wilson 算子的联合特征膜 & \S6.4 p.97 & 在 $\Psi=\infty$ 一侧 \\
\addlinespace[1pt]
magnetic eigenbrane & 磁特征膜 & 't Hooft 算子的联合特征膜 & \S6.4 p.97 & 在 $\Psi=0$ 一侧 \\
\addlinespace[1pt]
Hecke eigensheaf & Hecke 特征层 & Hecke 算子的本征对象，通常是 $\Dmod$-模 & \S6.4 p.95 / \S10.4 p.169 & 不是本征函数 \\
\addlinespace[1pt]
Hecke operator & Hecke 算子 & 由 Hecke 对应诱导、作用在丛或层上的函子 & \S9.1 p.132 & 不是普通的积分算子 \\
\addlinespace[1pt]
Hecke modification & Hecke 修正 & 在 $p_0$ 处允许极点、$p_0$ 外保持同构的重写 & \S9.1 p.132 / \S9.3 p.145 & $p_0$ 之外完全不改 \\
\addlinespace[1pt]
Hecke correspondence & Hecke 对应 & $Q(\Ldual w,z)\to\mathrm{Bun}_G\times\mathrm{Bun}_G$ & \S9.3 p.150 & 纤维中的点才是修正后的丛 \\

\addlinespace[4pt]
\multicolumn{5}{@{}l}{\tabhead{F\quad 单极子、Bogomolny 方程与仿射 Grassmann（\S9--\S10）}}\\
\addlinespace[1pt]

monopole & 单极子 & 't Hooft 算子附近的奇异 Dirac 解 & \S6.2 p.82 & 不是四维中可观测的粒子 \\
\addlinespace[1pt]
monopole bubbling & 单极子冒泡 & 瞬子缩成一点并并入已有单极子的过程 & \S10.1 p.148 / \S10.2 p.155 & 不是量子隧穿，是尺度模的经典退化 \\
\addlinespace[1pt]
Bogomolny equations & Bogomolny 方程 & $F=\ii\star D\varphi_0$ 的自对偶方程 & \S9 p.124 & 是 Hitchin 方程在 $\RR\times C$ 上的版本 \\
\addlinespace[1pt]
affine Grassmannian & 仿射 Grassmann & $G_{\CC}(X)/G_{\CC}(O)$ & \S9.1 p.132 / \S9.3 p.145 & 原文又称 loop Grassmannian \\
\addlinespace[1pt]
Schubert cell & Schubert 胞腔 & Hecke 修正空间上的轨道分层 & \S9.3 p.147 & 胞腔本身光滑，闭包一般奇异 \\
\addlinespace[1pt]
Schubert cycle & Schubert 循环 & Schubert 胞腔的闭包 & \S9.3 p.148 & 会分解为更低维的胞腔 \\
\addlinespace[1pt]
universal bundle & 泛丛 & $\MH\times C$ 上的丛，给出 $\xi(E)$ & \S7.2 p.106 & 不是 $\MH$ 上的丛 \\
\addlinespace[1pt]
hypermultiplet & 超多重态 & 四维 $N=2$ 的超多重态 & \S4.1 p.46 & 原文只在构造模空间时提一句 \\

\addlinespace[4pt]
\multicolumn{5}{@{}l}{\tabhead{G\quad Langlands 对应、D-模与几何背景（\S1、\S11）}}\\
\addlinespace[1pt]

geometric Langlands & 几何 Langlands & $\Ldual$ 的局部系统与 $\MC$ 上 Hecke 特征 $\Dmod$-模的对应 & \S1 p.2 & 不是数论 Langlands 的逐字搬运 \\
\addlinespace[1pt]
Langlands program & Langlands 纲领 & 数论与表示论之间的对偶纲领 & \S1 p.2 & 此处只借用其「对偶群」结构 \\
\addlinespace[1pt]
D-module & D-模 & 微分算子层 $\Dmod$ 的模 & \S11 p.184 & 不是「量子化之后的代数」 \\
\addlinespace[1pt]
sheaf of differential operators & 微分算子层 & 由 $\dd$ 与函数乘法生成的层 $\Dmod$ & \S1 p.5 / \S11.2 p.196 & 此处用的是扭化版本 $\Dmod_{K^{1/2}}$ \\
\addlinespace[1pt]
twisted D-module & 扭化 D-模 & 取值于 $K_M^{1/2}$ 的微分算子层的模 & \S11.2 p.196 & 与普通 $\Dmod$-模差一个线丛 \\
\addlinespace[1pt]
flat connection & 平坦联络 & 曲率为零的联络，等价于局部系统 & \S3.2 p.25 & 与 $G_{\CC}$-局部系统是同一份数据 \\
\addlinespace[1pt]
local system & 局部系统 & $\pi_1(C)$ 的表示，或局部常值层 & \S10.4 p.171（原文用 locally constant） & 原文没有用 local system 这个词 \\
\addlinespace[1pt]
fundamental group & 基本群 $\pi_1(C)$ & 闭曲线的同伦类构成的群 & \S5.6 p.77 & 对应的对象是 $\mathrm{Hom}(\pi_1(C),\Ldual)$ \\
\addlinespace[1pt]
irreducible representation & 不可约表示 & 没有非平凡子模的表示 & \S5.4 p.73 & 不必是酉表示 \\
\addlinespace[1pt]
ramified / unramified & 分歧 / 非分歧 & $C$ 上是否有分支点 & \S1 p.5 & 本文只处理非分歧情形 \\
\addlinespace[1pt]
principal bundle & 主丛 & 纤维为 $G$ 的丛 & \S7.1 p.108 & 特征类 $\xi(E)$ 是主丛的特征类 \\
\addlinespace[1pt]
vector bundle & 向量丛 & 纤维为向量空间的丛 & \S3.3 p.30 & 本文的 $E$ 默认是向量丛或主丛 \\
\addlinespace[1pt]
moduli space & 模空间 & 对象按同构分类所得的空间 & \S1 p.3 & 必须加稳定性才良定义 \\
\addlinespace[1pt]
Casimir polynomial & Casimir 多项式 & $\Tr\varphi^k$ 这类不变多项式 & \S4.3 p.53 & 独立 Casimir 的次数为 $d_i$ \\
\addlinespace[1pt]
quantization & 量子化 & 由经典系统构造量子理论 & \S1 p.4 & 与「取 D-模」不是同义词 \\
\addlinespace[1pt]
noncommutative geometry & 非交换几何 & 以非交换代数替代函数环的几何 & \S11 p.184 & 此处指余切丛的非交换形变 \\
\addlinespace[1pt]
Fourier--Mukai transform & Fourier--Mukai 变换 & 由核层给出的导出范畴等价 & \S1 p.4 & 此处是 T-对偶的数学化 \\
\addlinespace[1pt]
Kähler quotient & Kähler 商 & 一个 moment map 取零再取商 & \S4.1 p.40 & 超凯勒商要三个 moment map 同时取零 \\
\addlinespace[1pt]
flux & 通量 & 特征类 $\xi(E)$，或离散电、磁通量 $e,m$ & \S7.1 p.103 & 不是电磁场的通量密度 \\
\addlinespace[1pt]
de Rham differential & de Rham 微分 & 流形上的外微分 $\delta$ & \S4.1 p.45 & 本文的 $\delta$ 就是它 \\
\addlinespace[1pt]
vanishing theorem & 消灭定理 & 某些层上同调为零的定理 & \S5.1 p.65 & 用来说明 A-模型态空间是常值映射的上同调 \\

\addlinespace[4pt]
\multicolumn{5}{@{}l}{\tabhead{H\quad 原文未真正使用的背景术语（$\dagger$ 标注）}}\\
\addlinespace[1pt]

Nahm equations$^\dagger$ & Nahm 方程 & 单极子模空间满足的矩阵方程 & 全文未出现（仅引 Nahm 定理） & 原文没有给出方程本身 \\
\addlinespace[1pt]
Nahm's theorem$^\dagger$ & Nahm 定理 & 超对称杨--米尔斯理论所能存在的最大维数是十维 & \S2.1 p.6 & 与 Nahm 方程不是同一件事 \\
\addlinespace[1pt]
coherent sheaf$^\dagger$ & 相干层 & 局部有限表现的 $\mathcal{O}$-模 & 全文未出现 & 原文用层与复形，但不使用此术语 \\
\addlinespace[1pt]
derived category$^\dagger$ & 导出范畴 & 复形模去拟同构所得的范畴 & 全文未出现 & 本文的 Fourier--Mukai 是启发式用法 \\
\addlinespace[1pt]
loop group$^\dagger$ & 环路群 & $G_{\CC}((z))$ 这类群 & 全文未出现 & 原文只提 loop Grassmannian \\
\addlinespace[1pt]
opers$^\dagger$ & opers & 带额外结构的平坦联络 & 全文未出现 & 本文不涉及这种技术 \\
\addlinespace[1pt]
Dolbeault cohomology$^\dagger$ & Dolbeault 上同调 & $\bar\partial$ 的上同调 & 全文未出现 & 原文直接写 $H^q(X,\wedge^pT_X)$ \\
\addlinespace[1pt]
Kac--Moody algebra$^\dagger$ & Kac--Moody 代数 & 仿射李代数 & 全文未出现 & 原文用仿射 Grassmann 但不提它 \\

\end{longtable}
\endgroup

\noindent 表中有两类条目值得单独说明。第一类是\emph{原文一次也没有以该词形出现}的术语（标 $\dagger$ 的 Nahm 方程、相干层、导出范畴等）：它们常出现在二手文献和综述里，读者容易误以为本文用了，其实没有；表中给出的是原文的实际说法，例如用 \emph{flat connection} 或 \emph{locally constant} 代替 local system，用 $H^q(X,\wedge^pT_X)$ 代替 Dolbeault 上同调。第二类是\emph{同一个对象有三个名字}的术语：标量超荷、拓扑超荷、BRST 算子指的是同一个 $Q$\anchor{src p.18}，表中不合并成一条，是为了让三种叫法都能被检索到。

\subsection{符号索引}\label{sec:terms-symbols}

下面的表汇总正文与本章使用的记号。原则是：\emph{符号按原文的用法固定，不按数学文献的通行用法改}。例如 $\Psi$ 在原文是 canonical parameter 而不是某个波函数；$t$ 是扭化参数而不是时间；$B$ 既是 Hitchin 基又是 B-场，靠上下文区分，表中分别列出。\anchor{src p.35}

\begingroup
\footnotesize
\setlength{\tabcolsep}{5pt}
\renewcommand{\arraystretch}{1.1}
\begin{longtable}{@{}>{\raggedright\arraybackslash}p{3.15cm}>{\raggedright\arraybackslash}p{9.05cm}>{\raggedright\arraybackslash}p{3.1cm}@{}}
\toprule
\tabhead{符号} & \tabhead{含义} & \tabhead{首现} \\
\midrule
\endfirsthead
\toprule
\tabhead{符号} & \tabhead{含义} & \tabhead{首现} \\
\midrule
\endhead
\bottomrule
\endlastfoot

\multicolumn{3}{@{}l}{\tabhead{群、格与耦合常数}}\\
\addlinespace[1pt]
$G$ & 紧半单规范群 & \S1 p.2 \\
\addlinespace[1pt]
$\Ldual$ & Langlands 对偶群，由根与余根数据互换得到 & \S1 p.3 \\
\addlinespace[1pt]
$\mathfrak{g}$, $\gC$ & $G$ 的李代数及其复化 & \S3.1 p.19 \\
\addlinespace[1pt]
$G_{\CC}$, $G_{ad}$, $\tilde G$ & 复化群、伴随群、万有覆盖 & \S7 p.103 \\
\addlinespace[1pt]
$Z=\pi_1(G_{ad})$ & 中心，与自身有完美配对 $\Upsilon$ & \S7 p.103 \\
\addlinespace[1pt]
$\Lambda$, $\Lambda^\vee$ & 根格与余根格（电、磁荷的取值格） & \S7 p.104 \\
\addlinespace[1pt]
$\Upsilon$ & 中心的完美配对 $Z\times Z\to U(1)$ & \S7 p.104 \\
\addlinespace[1pt]
$\tau$ & 复耦合 $\theta/2\pi+4\pi\ii/e^2$ & \S2.2 p.11 \\
\addlinespace[1pt]
$\theta$, $e$ & theta 角与规范耦合 & \S2.2 p.11 \\
\addlinespace[1pt]
$n_g$ & 非单连通时的修正因子，取 $1,2,3$ & \S2.2 p.11 \\
\addlinespace[1pt]
$\Gamma$ & 对偶群，$SL(2,\ZZ)$ 或其 Hecke 子群 & \S2.2 p.12 \\

\addlinespace[3pt]
\multicolumn{3}{@{}l}{\tabhead{扭化与参数}}\\
\addlinespace[1pt]
$\kappa$ & 扭化所用的同态 $Spin(4)\to Spin(6)$ & \S3.1 p.17 \\
\addlinespace[1pt]
$t=v/u$ & 扭化参数，取值于 $\PP^1$；$t=\pm\ii$ 与 $t=\pm1$ 最重要 & \S3.2 p.25 \\
\addlinespace[1pt]
$\Psi$ & canonical parameter，式 (3.50)；理论真正依赖的参数 & \S3.5 p.35 \\
\addlinespace[1pt]
$u,v$ & 超荷 $Q=uQ_\ell+vQ_r$ 的系数 & \S3.2 p.25 \\
\addlinespace[1pt]
$\PP^1_g$, $\PP^1_h$ & 四维 TQFT 族的参数空间；$W$ 上的复结构族参数空间 & \S3.1 p.17 / \S4.1 p.43 \\

\addlinespace[3pt]
\multicolumn{3}{@{}l}{\tabhead{几何：$C$、$M$ 与 $\MH$}}\\
\addlinespace[1pt]
$C$, $g$ & 亏格 $g$ 的闭 Riemann 面；$g$ 为亏格 & \S1 p.2 \\
\addlinespace[1pt]
$M$ & 四维流形，扭化理论的底流形 & \S3 p.16 \\
\addlinespace[1pt]
$\Sigma$ & 二维世界面 & \S5 p.62 \\
\addlinespace[1pt]
$\MH$ & Hitchin 模空间，稳定 Higgs 丛的模空间 & \S4.1 p.45 \\
\addlinespace[1pt]
$\MC$ & 稳定 $G$-丛的模空间 & \S11 p.184 \\
\addlinespace[1pt]
$W$ & 全体 $(A,\varphi)$ 构成的无穷维仿射空间 & \S4.1 p.41 \\
\addlinespace[1pt]
$I,J,K$ & 三个复结构，满足 $I^2=J^2=K^2=IJK=-1$ & \S4.1 p.42 \\
\addlinespace[1pt]
$I_w$ & $w\in\PP^1_h$ 参数化的复结构族 & \S4.1 p.43 \\
\addlinespace[1pt]
$\omega_I,\omega_J,\omega_K$ & 三个辛形式；$\omega_J,\omega_K$ 恰当 & \S4.1 p.44 \\
\addlinespace[1pt]
$\mu$ & moment map，由 $\iota_{V(\epsilon)}\omega=\delta\mu$ 定义 & \S4.1 p.45 \\
\addlinespace[1pt]
$\nu_I=\mu_J+\ii\mu_K$ & 复结构 $I$ 下的全纯 moment map & \S4.2 p.49 \\
\addlinespace[1pt]
$B$, $d_i$ & Hitchin 基 $\bigoplus_iH^0(C,K_C^{d_i})$ 与 Casimir 次数 & \S4.3 p.53 \\
\addlinespace[1pt]
$\pi$ & Hitchin 纤维化 $\MH\to B$ & \S4.3 p.53 \\
\addlinespace[1pt]
$L$ & 行列式线丛，生成 $\mathrm{Pic}(\MH)\cong\ZZ$ & \S4.1 p.45 \\
\addlinespace[1pt]
$K_C$, $K_M$ & $C$ 与靶空间的典范丛 & \S4.2 p.49 / \S11.2 p.196 \\
\addlinespace[1pt]
$\Omega_J$ & $\omega_K+\ii\omega_I$，复结构 $J$ 的全纯辛形式 & \S4.1 p.44 \\

\addlinespace[3pt]
\multicolumn{3}{@{}l}{\tabhead{场与算子}}\\
\addlinespace[1pt]
$A$, $F$, $D$ & 联络、曲率、协变导数；原文约定 $D=\dd+A$ & \S3.1 p.19 / \S6.1 p.79 \\
\addlinespace[1pt]
$\varphi$ & Higgs 场，$K_C\otimes\mathrm{ad}E$ 的全纯截面 & \S4.2 p.49 \\
\addlinespace[1pt]
$\varphi_{0},\varphi_{1,2,3}$ & 扭化后留下的标量与 adjoint 值 $1$-形式 & \S3.1 p.19 \\
\addlinespace[1pt]
$A+\ii\varphi$ & 复联络；$t=1$ 时 Bogomolny 方程等价于 $F+\ii\star F=0$ & \S9 p.124 \\
\addlinespace[1pt]
$Q$ & 标量超荷，$Q^2=0$，定义 $Q$-上同调 & \S3.1 p.18 \\
\addlinespace[1pt]
$E$, $E(R)$ & 向量丛；表示 $R$ 的伴随丛 $E(\mathcal R)$ & \S4.2 p.49 / \S6.1 p.80 \\
\addlinespace[1pt]
$\mathbb{A}$ & 联络空间（$W$ 的第一分量） & \S4.1 p.41 \\
\addlinespace[1pt]
$B_0$, $\omega_0$ & B-场与辛形式，广义复几何中的 B-场变换用 & \S5.2 p.66 \\
\addlinespace[1pt]
$B$（重载） & Hitchin 基、B-场、B-膜，三者靠上下文区分 & \S4.3 p.53 / \S4.1 p.46 \\

\addlinespace[3pt]
\multicolumn{3}{@{}l}{\tabhead{膜、算子与 Hecke 数据}}\\
\addlinespace[1pt]
$H_{B,B'}$ & 膜 $B$ 与 $B'$ 之间的态空间 & \S6.4 p.93 \\
\addlinespace[1pt]
$V$, $V_p$ & 特征膜的本征值空间；$p$ 变时按平坦丛纤维变 & \S6.4 p.95 \\
\addlinespace[1pt]
$q$ & Wilson 算子的电权，记 $E_q(R)$ & \S8.1 p.116 \\
\addlinespace[1pt]
$e_0$, $m_0$ & 二维紧化后的离散电、磁通量 & \S7.2 p.110 \\
\addlinespace[1pt]
$\xi(E)$ & 主丛的特征类，属 $H^2(M,\pi_1(G))$ & \S7 p.103 \\
\addlinespace[1pt]
$\zeta=\xi(E_{ad})$ & 泛丛的特征类，$\xi(E)=(\Phi\times1)^*\zeta$ & \S7.2 p.110 \\
\addlinespace[1pt]
$e$, $m$ & 离散电通量与磁通量 & \S7 p.104 \\
\addlinespace[1pt]
$\Ldual w$, $T(\Ldual w)$ & $\Ldual$ 的最高权与对应的 't Hooft 算子 & \S9.1 p.132 \\
\addlinespace[1pt]
$T_{(k)}$ & $\wedge^kV$ 对应的 't Hooft 算子，$k<N$ 时不可逆 & \S9.1 p.133 \\
\addlinespace[1pt]
$\Ldual R$ & 对偶群的表示，给 't Hooft 算子贴标签 & \S6.2 p.81 \\
\addlinespace[1pt]
$Y(\Ldual w)$, $Y(\rho)$ & Hecke 修正空间 / Schubert 胞腔 & \S9.3 p.147 \\
\addlinespace[1pt]
$N_\alpha$, $n_\alpha$ & OPE 的向量空间（开线）与整数（闭环） & \S10.4 p.170 / \S6.4 p.93 \\
\addlinespace[1pt]
$Z(\Ldual R_i,p_i)$ & 带 't Hooft 插入的 Bogomolny 解模空间 & \S10.4 p.170 \\

\addlinespace[3pt]
\multicolumn{3}{@{}l}{\tabhead{层与模空间}}\\
\addlinespace[1pt]
$\Dmod$ & 微分算子层 & \S1 p.5 \\
\addlinespace[1pt]
$\Dmod_{K^{1/2}}$ & 取值于 $K_M^{1/2}$ 的扭化微分算子层 & \S11.2 p.196 \\
\addlinespace[1pt]
$\mathrm{Bun}_G$ & 所有 $G$-丛构成的模 stack & \S9.3 p.150 \\
\addlinespace[1pt]
$Gr_N$, $Gr_G$ & 仿射 Grassmann 空间（$U(N)$ 与一般 $G$） & \S9.3 p.145 \\
\addlinespace[1pt]
$X$, $O$ & $C^*$ 上全纯函数环；$p_0$ 处全纯函数环 & \S9.3 p.145 \\
\addlinespace[1pt]
$\CC((z))$, $\CC[[z]]$ & 形式 Laurent 级数 / 形式幂级数 & \S9.3 p.145 \\

\end{longtable}
\endgroup

\noindent 与 \texttt{plan/02} 符号卡的差异只有两处，且都以原文为准：其一，符号卡把 $\tau$ 当作主角，而原文 \S3.5 引入 canonical parameter $\Psi$ 并强调理论只依赖 $\Psi$\anchor{src p.35}，$\tau$ 只在 $t=\pm\ii$ 时完全退出；其二，符号卡没有列 $\nu_I=\mu_J+\ii\mu_K$（全纯 moment map），而原文 \S4.2 正是用它把 Hitchin 方程写成 $\nu_I=0$。其余符号两处一致。若将来两处冲突，一律以原文的式 (3.50) 与 (4.20) 为准。\anchor{src p.49}

\subsection{常见误解}\label{sec:terms-misconceptions}

以下 25 条不是「翻译错误」而是「读法错误」：每一条的「错在哪」都尽量说明这个直觉是从哪个正确的邻域类推来的，为什么自然；「正确图像」都回指原文锚点或骨架节点（S1--S8、M1--M5），方便对照。

\paragraph{误解 1：扭化就是把 $N=4$ 理论简化成一个更小的理论}
\mbox{}\par\noindent
\textbf{错在哪}：扭化之后只保留 $Q$-闭的态、把 $[Q,\mathcal{O}]$ 当作零，看起来确实像丢掉了大半自由度；再加上最常见的 Donaldson--Witten 扭化是从 $N=2$ 里得到更小的拓扑理论，“扭化＝截断”的类比非常自然。
\par\noindent
\textbf{正确图像}：扭化是重定义自旋，不是截断：把 $Spin(4)$ 换成 $Spin(4)\times Spin(6)$ 中与之同构的 $Spin'(4)$，它在 $\RR^4$ 上的作用不变，但使 $N=4$ 理论的作用量以新的方式分解。理论还是同一个理论，改的是“什么算算子、什么算态”的约定。见 \S3.1\anchor{src p.17}，骨架 S2/N3。

\paragraph{误解 2：扭化后只留 $Q$-不变的态，所以理论变平凡了}
\mbox{}\par\noindent
\textbf{错在哪}：物理文献常说“只保留 $Q$-不变的东西”，这句省略了后半句；如果只取不变子空间，得到的仍然庞大且依赖度量，“拓扑”两字就无从谈起。
\par\noindent
\textbf{正确图像}：算子与态落在 $Q$ 的上同调群里：只留 $Q$-闭的，并把 $Q$-恰当的视为零。这样度量或耦合的连续形变才不改变结果；原文把“能量--动量张量是 $Q$-恰当的”作为拓扑场论的判据，并说这种理论只是“松散地”算 TQFT。见 \S3.1\anchor{src p.18}，骨架 S2。

\paragraph{误解 3：A-模型就是 A-膜，B-模型就是 B-膜}
\mbox{}\par\noindent
\textbf{错在哪}：两对词都带 A/B，而且经常出现在同一句话里（例如“$J$ 的 B-模型里的 A-膜”），中文里“模型”“膜”又只差一个字。
\par\noindent
\textbf{正确图像}：模型是一套二维场论（路径积分、拓扑超荷、可积的形变），膜是这套场论的边界条件。给定一个二维模型，膜的 A/B 型取决于它相对靶空间某个复结构是 A 型还是 B 型；同一个模型可以同时容纳不同类型的膜——本文 \S5.6 就在 $\MH$ 上同时讨论四种类型的膜。见 \S5.1\anchor{src p.65} 与 \S5.6\anchor{src p.77}。

\paragraph{误解 4：镜像对称就是两个 Calabi--Yau 流形几何等同}
\mbox{}\par\noindent
\textbf{错在哪}：镜像对称的经典例子是 CY 三流形，于是“镜像”被读成“同一个流形的另一种几何”，甚至以为两边存在自然的微分同胚可以逐点对照。
\par\noindent
\textbf{正确图像}：镜像对称是两个 sigma 模型的等价：一边的 A-模型等于另一边的 B-模型。本文的镜像对是 $\MH(G,C)$ 与 $\MH(\Ldual,C)$，它们一般不是 CY 流形，而且只在复结构 $J$ 下成镜像对；换一个复结构就不是镜像对了。见 \S5.3\anchor{src p.70}，骨架 S4/M2。

\paragraph{误解 5：镜像对的两侧是同一个流形，只是换了个复结构}
\mbox{}\par\noindent
\textbf{错在哪}：$\MH(G)$ 与 $\MH(\Ldual)$ 的定义几乎逐字对称，而且 S-对偶只换群，容易以为靶空间没变。
\par\noindent
\textbf{正确图像}：两侧维数相同（都是 $2\dim B$），但一般并不同构：$G=SU(N)$ 时对偶是 $PSU(N)$，只有自对偶的群（如 $E_8$）才可能两边相同。对偶性来自“换群 + 换量子化方式”两件事的合成，而不只是换复结构。见 \S5.3\anchor{src p.70}，M2。

\paragraph{误解 6：$\MH$ 就是内部紧化空间}
\mbox{}\par\noindent
\textbf{错在哪}：紧化后出现的模空间容易被当成“额外维度的形状”；$\MH$ 又确实是某个几何对象的模空间，名字里也带“模空间”，读起来很像在描述内部流形。
\par\noindent
\textbf{正确图像}：$\MH$ 是二维低能有效理论（sigma 模型）的靶空间，是四维理论紧化到 $C$ 上之后\emph{涌现}出来的，不是被紧化掉的那部分空间的形状。原文的紧化方向是 $4d\to2d$，被紧化的是 $C$ 之外的两个方向。见 \S1\anchor{src p.3} 与 \S4\anchor{src p.40}，骨架 S3/N4。

\paragraph{误解 7：S-对偶就是 T-对偶}
\mbox{}\par\noindent
\textbf{错在哪}：两者都反演一个参数（$\tau\mapsto-1/\tau$ 与半径 $R\mapsto1/R$），都能写成某个群的生成元，中文又都译作“对偶”，在二维时甚至能出现在同一个公式里。
\par\noindent
\textbf{正确图像}：S-对偶是四维理论内部的量子对称：换 $G\leftrightarrow\Ldual$、换电与磁、$\tau\mapsto-1/n_g\tau$；T-对偶是紧化到二维后才出现的现象：沿 sigma 模型的环面纤维反演半径，等价于对 $[B+\ii\omega]$ 作平移。本文的关键链条恰恰是“四维的 S-对偶 = 二维的镜像对称”，而镜像对称再实现为纤维上的 T-对偶。见 \S1\anchor{src p.3} 与 \S5.5\anchor{src p.77}，S4/M2。

\paragraph{误解 8：S-对偶只是把 $G$ 换成 $\Ldual$}
\mbox{}\par\noindent
\textbf{错在哪}：$G\leftrightarrow\Ldual$ 是最显眼、最容易在表里查到的一半（Table 1 只列了群对），而且符号上 $\tau$ 的变换看起来只是一个 Möbius 变换。
\par\noindent
\textbf{正确图像}：S-对偶同时交换电与磁的源（Wilson 与 't Hooft）、交换 $\Psi\mapsto-1/n_g\Psi$，并且在对 $\MH$ 上的膜作用时交换两种量子化（A-模型与 B-模型）。只换群而漏掉“换量子化”，对膜的作用就会算错方向。见 \S2.2\anchor{src p.11} 与 \S7.2\anchor{src p.116}，骨架 S1。

\paragraph{误解 9：canonical parameter 就是 $\tau$，或者只是把 $\tau$ 换个记号}
\mbox{}\par\noindent
\textbf{错在哪}：\S3.1--\S3.4 一直在用 $\tau$，\S3.5 突然改用 $\Psi$，最自然的猜测是“$\Psi$ 是 $\tau$ 的某个函数”；而 $\tau$ 的定义早在 \S2.2 就给过了，看起来没有引入新参数的必要。
\par\noindent
\textbf{正确图像}：$\Psi$ 由作用量中 $I_2$ 项的系数定义（式 (3.50)），它把 $\tau$ 与扭化参数 $t$ 混在一起；理论只依赖 $\Psi$，而 $\tau$ 只在 $t=\pm\ii$ 时完全退出（此时 $\Psi=\infty$）。这才是“几何 Langlands 真正相关的参数不是 $\tau$”的意思，也解释了为什么 T-对偶（$\Psi\mapsto\Psi+1$）与 S-对偶（$\Psi\mapsto-1/n_g\Psi$）都表现为 Möbius 变换。见 \S3.5\anchor{src p.35}，骨架 S2/M1。

\paragraph{误解 10：'t Hooft 算子就是 Wilson 圈的对偶}
\mbox{}\par\noindent
\textbf{错在哪}：原文说 S-对偶把 $\Psi=\infty$ 处的 Wilson 算子送到 $\Psi=0$ 处的 't Hooft 算子，“互为对偶”是这句话最顺口的改写；而且两者都记作沿一条线的算子。
\par\noindent
\textbf{正确图像}：更准确的说法是：S-对偶是参数空间上的对称，它把“某个参数值处的一族算子”对应到“另一个参数值处的一族算子”。在同一个理论里，Wilson 是 order operator（把经典量再量子化），'t Hooft 是 disorder operator（规定经典场的奇异性再作路径积分），二者的构造方式不同，不是同一个算子的两个名字；只有在 S-对偶联系的两个参数点之间才有一一对应。见 \S6.2\anchor{src p.81}，骨架 S5。

\paragraph{误解 11：Hecke 特征层就是某个函数空间里的本征函数}
\mbox{}\par\noindent
\textbf{错在哪}：\emph{eigensheaf} 里含 eigen，Hecke 算子又被类比成算子，于是被读成“本征函数”；中文“特征层”里的“特征”也不提示区别。
\par\noindent
\textbf{正确图像}：本征对象是层（最终是 D-模），“本征值”是一个向量空间：$XB=B\otimes V$。移动参数 $p$ 时 $V$ 按 $C$ 上平坦丛的纤维变化，所以“本征值”有单值性，但没有函数可言。原文 \S6.4 定义 eigenbrane 时明确与本征值空间对照，\S10.4 再把它与 Hecke 特征层对应起来。见 \S6.4\anchor{src p.95}，骨架 S7/M3。

\paragraph{误解 12：Hecke 算子是普通的积分算子}
\mbox{}\par\noindent
\textbf{错在哪}：Hecke 算子常被写成对权格求和，形式上像一个平均积分；在 $SL(2,\ZZ)$ 语境里“Hecke 算子”也确实是一种平均。
\par\noindent
\textbf{正确图像}：本文的 Hecke 算子是 Hecke 对应 $Q(\Ldual w,z)\subset\mathrm{Bun}_G\times\mathrm{Bun}_G$ 诱导的函子：先拉回到对应空间，再推回 $\mathrm{Bun}_G$，作用在丛或层上而不是函数上。$U(N)$ 情形下 't Hooft 算子的作用就是 $L_i\mapsto L_i\otimes\mathcal{O}(p_0)^{m_i}$，与数学文献中 Hecke 算子的描述逐字一致。见 \S9.1\anchor{src p.132}，骨架 S7。

\paragraph{误解 13：D-模就是量子化}
\mbox{}\par\noindent
\textbf{错在哪}：D-模由微分算子生成，“微分算子＝动量”的量子化直觉很容易把它读成“把经典系统量子化”；非交换几何也常与量子化同义使用，两者叠在一起更强化了这个印象。
\par\noindent
\textbf{正确图像}：D-模是“微分算子层的模”这一层纯代数对象。本文用它，是因为 $\MC$ 上的 D-模可以看成其余切丛的非交换形变上的 B-膜，而 A-膜范畴与该非交换形变上的 B-膜范畴等价——这是补齐几何 Langlands 缺环的技术内容，而不是对某个经典系统作量子化的操作。见 \S11\anchor{src p.184}，骨架 S8。

\paragraph{误解 14：eigenbrane 就是 Hilbert 空间里的本征态}
\mbox{}\par\noindent
\textbf{错在哪}：eigen- 前缀加上“算子作用于对象”的句式，与量子力学里的本征态几乎同构；中文“特征膜”也不提示区别，很容易把它想成一个波函数。
\par\noindent
\textbf{正确图像}：eigenbrane 是被线算子映到自身的膜：$XB=B\otimes V$。膜不是 Hilbert 空间中的向量，而是二维场论的边界条件；“本征值”是 Chan-Paton 空间（向量空间或层），不是数。原文特意强调这一点，并指出它正是 Hecke 特征层概念的物理版本。见 \S6.4\anchor{src p.95}，骨架 S6。

\paragraph{误解 15：几何 Langlands 只是数论 Langlands 的逐字搬运}
\mbox{}\par\noindent
\textbf{错在哪}：两边都叫 Langlands，右端都出现“对偶群、自守对象”这类词，最省力的理解就是把数域换成曲线、其余照抄。
\par\noindent
\textbf{正确图像}：几何版本对应的是 $\Ldual$ 的局部系统与 $\MC$ 上“同时是 D-模的 Hecke 特征层”之间的对应；数论情形里的自守表示在这里的类比物是 Hecke 特征层，而不是函数。函数域情形的证明并不能直接给出几何版本，几何版本的关键工具是 Hitchin 系统与 S-对偶。原文明说：Hitchin 系统只是几何 Langlands 的半经典近似。见 \S1\anchor{src p.4} 与 \S11\anchor{src p.184}，S8。

\paragraph{误解 16：几何 Langlands 的右边是 $\MC$ 上的函数}
\mbox{}\par\noindent
\textbf{错在哪}：就像数论情形把表示实现成自守形式一样，人们期待“对应”的右边是某个函数空间；而左边 $\mathrm{Hom}(\pi_1(C),\Ldual)$ 也确实是参数空间，右边像是它的函数。
\par\noindent
\textbf{正确图像}：右边是 $\MC$ 上的 D-模（Hecke 特征层），也就是某个非交换形变上的 B-膜；不是 $\MH$ 上的函数，也不是 $\MC$ 上的函数。函数只出现在半经典层面：$\MH$ 在复结构 $I$ 下的全纯函数就是 Hitchin 基 $B$ 上的全纯函数，S-对偶把 $B$ 送到 ${}^LB$。见 \S11\anchor{src p.184} 与 \S5.4\anchor{src p.72}，S8。

\paragraph{误解 17：典范余各向同性膜就是某个具体的 Lagrange 膜}
\mbox{}\par\noindent
\textbf{错在哪}：名字里有“典范”，容易联想到余切丛的零截面之类的标准子流形；而“余各向同性”与“Lagrange”在辛几何教材里常常相邻出现，便以为是同一回事。
\par\noindent
\textbf{正确图像}：它是秩 $1$ 的 A-膜，支撑在\emph{整个} $\MH$ 上，Chan-Paton 是 $U(1)$ 线丛，条件 $(\omega^{-1}F)^2=-1$ 使 $\omega-\ii F$ 成为全纯辛形式。它不是 Lagrange 子流形（维数等于环境维数），而是“余各向同性且达到最大维数”的极端情形；它的镜像仍然是 $(A,B,A)$ 型。见 \S11.1\anchor{src p.185}，骨架 M5。

\paragraph{误解 18：膜的类型 $(B,B,B)$ 表示“三种 B-膜”}
\mbox{}\par\noindent
\textbf{错在哪}：三个字母排成一列，很像三个分量或三种并列的选择；“(B,A,A)”读起来也像“一个 B 加两个 A”的清单。
\par\noindent
\textbf{正确图像}：三个位置分别对应三个复结构 $I,J,K$，每一位写该膜相对那个复结构是 A 型还是 B 型。本文用到四种：$(B,B,B)$、$(B,A,A)$、$(A,B,A)$、$(A,A,B)$。所以 $(B,B,B)$ 是“对三个复结构都是 B 型”的\emph{一种}膜，不是三种膜。见 \S5.6\anchor{src p.77}。

\paragraph{误解 19：zerobrane 是零维的膜，所以没什么内容}
\mbox{}\par\noindent
\textbf{错在哪}：名字带 zero，支撑在一点，看起来像一个退化的平凡对象；而“零维”在拓扑里也常意味着信息很少。
\par\noindent
\textbf{正确图像}：zerobrane 是最重要的探针：它的 Chan-Paton 是支撑在该点的天穹层，被 Wilson 算子作用后给出 $E_q(R)|_x$，这直接说明它是电特征膜、本征值是向量空间。S-对偶把它送到 Hitchin 纤维上、带平坦 Chan-Paton 丛的膜（type F），这正是磁特征膜的原型。见 \S5.5\anchor{src p.74} 与 \S8.2\anchor{src p.121}，S6/M4。

\paragraph{误解 20：Hitchin 纤维化就是“忘掉 $\varphi$”的那个纤维化}
\mbox{}\par\noindent
\textbf{错在哪}：两个映射都从 $\MH$ 出发、都“丢掉信息”，而“第二纤维化”忘掉 $\varphi$ 的说法在文献里同样常见，很容易张冠李戴。
\par\noindent
\textbf{正确图像}：第二纤维化忘掉 $\varphi$、只留稳定 $G$-丛，靶是 $\mathcal{M}$，维数是 $\frac12\dim\MH$；Hitchin 纤维化 $\pi:\MH\to B$ 保留 $\varphi$ 的 Casimir 不变量 $\Tr\varphi^{d_i}$，靶是 $B=\bigoplus_iH^0(C,K_C^{d_i})$，典型纤维是阿贝尔簇（环面）。可积系统的环面纤维来自后者，不是前者。见 \S4.3\anchor{src p.52}。

\paragraph{误解 21：可积系统就是能显式解出来}
\mbox{}\par\noindent
\textbf{错在哪}：Hitchin 系统被称为“完全可积”，而经典力学里完全可积的例子（谐振子、开普勒问题）往往能写出解，于是以为每个纤维都能参数化。
\par\noindent
\textbf{正确图像}：完全可积只保证存在 $\frac12\dim\MH$ 个 Poisson 交换的 Hamilton 量（$B$ 上的线性函数），使 $\MH$ 分解为环面纤维上的 Liouville 结构；具体纤维是紧化 Prym 型的阿贝尔簇，一般不能显式解出，“可积”是关于结构的存在性而非求解公式。见 \S4.3\anchor{src p.54}。

\paragraph{误解 22：广义复几何就是复几何与辛几何的并集}
\mbox{}\par\noindent
\textbf{错在哪}：两个基本例子 $I_{\hat J}$ 与 $I_\omega$ 分别对应 B-模型与 A-模型，于是被读成“把两种几何放在一起”，像一个大伞下的两个分支。
\par\noindent
\textbf{正确图像}：它是在 $T\oplus T^*$ 上、带自然不定度量的结构理论：广义近复结构满足 $I^2=-1$、保度量、可积。$I_{\hat J}$ 与 $I_\omega$ 只是两个极端情形；B-场变换 $I(B_0)=M(B_0)IM(B_0)^{-1}$ 把一个连续地变成另一个，镜像对称才有可能在广义复几何的框架里被表述。见 \S5.2\anchor{src p.65}，骨架 S4。

\paragraph{误解 23：超凯勒商就是 Kähler 商}
\mbox{}\par\noindent
\textbf{错在哪}：Hitchin 方程被写成“三个 moment map 为零”，看起来像把 Kähler 商重复三次；而“商”的构造在两种情形下确实都是取零集再除群。
\par\noindent
\textbf{正确图像}：Kähler 商只用一个 moment map（一个辛形式），得到的流形带一个复结构；超凯勒商要求 $\mu_I,\mu_J,\mu_K$ 三个 moment map 同时为零，再取商，得到的流形同时带三个复结构。若只取一个 moment map 为零，得到的只是其中一个复结构的商。见 \S4.1\anchor{src p.40}。

\paragraph{误解 24：稳定性条件只是技术性假设，可以忽略}
\mbox{}\par\noindent
\textbf{错在哪}：原文定义 $\MH$ 时把“稳定”压缩在一句话里，之后的计算几乎只用 Hitchin 方程，稳定性看起来只是保证模空间存在的技术细节。
\par\noindent
\textbf{正确图像}：不稳定解必须丢掉，否则模空间不对、Hitchin 纤维化与可积性都不成立。\S3.2 的脚注已经说明 $Y_M(G)$ 需要去掉不稳定轨道，\S4.2 又给了一般 $G$ 的稳定性定义（用极大抛物子群约化），并说明 $\MH$ 是稳定 Hitchin 对的模空间。见 \S4.2\anchor{src p.50}。

\paragraph{误解 25：单极子冒泡是量子隧穿之类的量子过程}
\mbox{}\par\noindent
\textbf{错在哪}：bubbling 在物理里常指真空泡或量子跃迁，“瞬子缩成一点”听起来也像量子涨落；而单极子本身在本文又是 't Hooft 算子附近的经典奇异解，更容易让人想到量子修正。
\par\noindent
\textbf{正确图像}：冒泡是经典尺度模的退化：Bogomolny 解在 $C^2\cong\RR^4$ 上对应等变瞬子，瞬子方程具有尺度不变性，所以瞬子可以缩成一点并“并入”已有的单极子；数学上，这对应 \S9.3 中 Hecke 修正空间闭包的分层——闭包中出现更低维的轨道，就是“冒泡”的几何影子。见 \S10.2\anchor{src p.155}，骨架 S7。

\vspace{1ex}
\noindent 这 25 条的共同结构是：一个在别处正确、在此处稍有错位的直觉。真正的读法不是记住结论，而是记住“这个直觉的邻域在哪里失效”——多半失效在三个地方：把二维与四维的目标搞混（误解 6、7）、把“模型”与“膜”搞混（误解 3、4、18）、把“函数式”的语言硬套到层上（误解 11、12、14、16）。这三类混淆正好对应本章三张表的分工：术语表管第一类，符号索引管第二类，本节的“正确图像”管第三类。

